\section{Síntesis y ampliación}

\subsection{Una cadena de razonamiento que atraviesa toda la clase}

Esta lección se puede resumir en seis pasos. Primero, lleva un registro por separado para los parámetros, el estado del optimizador, las gradientes, las activaciones y los flujos de datos. Segundo, realiza el particionamiento (sharding) a lo largo del eje lógico más coherente. Tercero, escribe las dependencias colectivas o punto a punto necesarias para restaurar la semántica de un solo dispositivo. Cuarto, determina si la comunicación se encuentra en la ruta crítica o puede superponerse. Quinto, mapea los ejes lógicos a la topología física real: NVLink, NVSwitch, InfiniBand, PCIe y almacenamiento. Sexto, valida mediante el profiler, el rendimiento por unidad de tiempo (throughput), la memoria pico, la utilización de red y las curvas de convergencia, en lugar de aceptar basándose únicamente en los nombres de configuración.

\begin{importantbox}{El núcleo de Ultra-scale no son más GPUs, sino menos esperas innecesarias}
Aumentar el número de dispositivos incrementa solo la paralelidad potencial. La paralelidad potencial solo se convierte en rendimiento efectivo de entrenamiento cuando el estado está correctamente particionado, los colectivos están correctamente agrupados, el micro-lote (micro-batch) llena el cronograma (schedule), la carga de ruta se equilibra y los datos y los puntos de control (checkpoint) no provocan que el cómputo quede sin trabajo.
\end{importantbox}

\subsection{Lista de verificación práctica}

\begin{enumerate}
  \item Registra el tipo de dato (dtype), la forma (shape), el factor de replicación y el ciclo de vida para cada categoría de estado.
  \item Marca para cada eje paralelo el grupo de procesos, el enlace físico y los colectivos críticos.
  \item Mide por separado el tiempo del paso solo computacional, solo comunicacional y de extremo a extremo.
  \item Verifica si la superposición reduce el tiempo total, en lugar de solo hacer que las líneas de tiempo se solapen.
  \item Registra los retrasados (stragglers), el desequilibrio del router, las burbujas del pipeline y las paradas por entrada o checkpoint.
  \item Replica la pérdida (loss) con una sola tarjeta a la escala mínima, luego incrementa la paralelidad eje por eje y compara los valores numéricos y el rendimiento.
  \item Guarda un punto de control (checkpoint) recuperable y realiza una prueba real del tiempo de recuperación tras un fallo de nodo.
\end{enumerate}

\subsection{Lecturas adicionales}

\begin{itemize}
  \item Hugging Face, \href{https://huggingface.co/spaces/nanotron/ultrascale-playbook}{Ultra-Scale Playbook}: el diagrama de sistema principal y el marco de toma de decisiones de esta lección.
  \item Rajbhandari et al., \href{https://arxiv.org/abs/1910.02054}{ZeRO}; PyTorch, \href{https://docs.pytorch.org/docs/stable/distributed.fsdp.fully_shard.html}{documentación de FSDP2 fully\_shard}.
  \item Liu et al., \href{https://arxiv.org/abs/2310.01889}{Ring Attention con Transformers por bloques}.
  \item DeepSeek-AI, \href{https://arxiv.org/abs/2412.19437}{Informe técnico de DeepSeek-V3} y \href{https://github.com/deepseek-ai/DeepEP}{DeepEP}.
  \item \href{https://github.com/pytorch/torchtitan}{TorchTitan}, \href{https://github.com/huggingface/nanotron}{Nanotron}, \href{https://github.com/NVIDIA/Megatron-LM}{Megatron-LM}: implementaciones ejecutables de paralelismo multidimensional.
  \item Hugging Face, \href{https://huggingface.co/spaces/HuggingFaceTB/smol-training-playbook}{Smol Training Playbook}; \href{https://jax-ml.github.io/scaling-book/tpus/}{JAX Scaling Book}: infraestructura de entrenamiento y perspectiva desde TPU.
\end{itemize}

\end{document}
